\section{MDL, chọn cấu trúc và streaming}
\label{sec:mdl-streaming}

\subsection{MDL bằng một hình ảnh rất đời thường}

Giả sử bạn phải gửi một file dữ liệu cho đồng nghiệp qua một kênh rất chậm. Bạn có hai cách:
gửi nguyên file, hoặc gửi một ``công thức mô tả file'' cộng với phần dữ liệu không khớp công thức.
Nếu công thức quá đơn giản, phần sai lệch còn lại rất lớn. Nếu công thức quá phức tạp, bản thân
công thức đã dài. \textbf{Minimum Description Length} (MDL) chọn cách mô tả có \emph{tổng số bit}
nhỏ nhất:
\begin{equation}
  \bigl\langle \mathcal{X};\mathcal{F}\bigr\rangle
  =\bigl\langle \mathcal{F}\bigr\rangle
   +\bigl\langle \mathcal{X}\mid\mathcal{F}\bigr\rangle.
  \label{eq:mdl-total}
\end{equation}
Đây là Eq.~(10) trong paper. Số hạng thứ nhất là chi phí mô tả mô hình; số hạng thứ hai là chi phí
mã hóa dữ liệu khi đã biết mô hình. Mô hình tốt không phải mô hình khớp tuyệt đối, mà là mô hình
giải thích được nhiều dữ liệu với ít bit.

\begin{keypoint}
MDL là một nguyên tắc cân bằng giữa \textbf{độ khớp} và \textbf{độ phức tạp}. Nó đồng thời dùng
để chọn số chiều ẩn, quyết định giữ outlier nào, và quyết định có thêm một regime mới hay không.
Vì vậy paper mới có thể nói mô hình ``tự động''.
\end{keypoint}

\subsection{Chi phí mô tả mô hình}

Với mỗi mô hình $\Theta$, paper tính chi phí từ các tham số khác 0. Gọi $|\cdot|$ là số phần tử
khác 0, $c_F=32$ là số bit cho một số thực, và $\log^*$ là độ dài mã phổ dụng cho một số nguyên.
Các công thức có cùng một khuôn:
\begin{align}
  \bigl\langle\mathcal{F}\bigr\rangle
  &=\sum_{\Theta\in\mathcal{F}}\bigl\langle\Theta\bigr\rangle,
  &
  \bigl\langle\Theta\bigr\rangle
  &=\bigl\langle\theta_d\bigr\rangle
   +\bigl\langle\theta_s\bigr\rangle
   +\bigl\langle\widehat{\mathcal{X}}_o\bigr\rangle,
  \label{eq:model-cost}\\
  \bigl\langle\mathbf{W}^{(\mathrm{key})}\bigr\rangle
  &=|\mathbf{W}^{(\mathrm{key})}|
    \bigl(\log d_k+\log k+c_F\bigr)
    +\log^*|\mathbf{W}^{(\mathrm{key})}|,
  \label{eq:cost-wkey}\\
  \bigl\langle\mathbf{W}^{(\mathrm{loc})}\bigr\rangle
  &=|\mathbf{W}^{(\mathrm{loc})}|
    \bigl(\log d_l+\log l+c_F\bigr)
    +\log^*|\mathbf{W}^{(\mathrm{loc})}|,\\
  \bigl\langle\theta\bigr\rangle
  &=\bigl\langle\mathbf{A}\bigr\rangle
   +\bigl\langle\mathcal{D}\bigr\rangle,\\
  \bigl\langle\mathbf{A}\bigr\rangle
  &=|\mathbf{A}|\bigl(\log d_k+\log d_l+c_F\bigr)
    +\log^*|\mathbf{A}|,\\
  \bigl\langle\mathcal{D}\bigr\rangle
  &=|\mathcal{D}|\bigl(\log d_k+2\log d_l+c_F\bigr)
    +\log^*|\mathcal{D}|.
  \label{eq:cost-theta}
\end{align}
Các hệ số mùa vụ và outlier có chi phí tương tự:
\begin{align}
  \bigl\langle\mathbf{S}^{(\mathrm{time})}\bigr\rangle
  &=|\mathbf{S}^{(\mathrm{time})}|
    \bigl(\log d_s+\log L_c+c_F\bigr)
    +\log^*|\mathbf{S}^{(\mathrm{time})}|,\\
  \bigl\langle\mathbf{S}^{(\mathrm{key})}\bigr\rangle
  &=|\mathbf{S}^{(\mathrm{key})}|
    \bigl(\log d_s+\log k+c_F\bigr)
    +\log^*|\mathbf{S}^{(\mathrm{key})}|,\\
  \bigl\langle\mathbf{S}^{(\mathrm{loc})}\bigr\rangle
  &=|\mathbf{S}^{(\mathrm{loc})}|
    \bigl(\log d_s+\log l+c_F\bigr)
    +\log^*|\mathbf{S}^{(\mathrm{loc})}|,\\
  \bigl\langle\widehat{\mathcal{X}}_o\bigr\rangle
  &=|\widehat{\mathcal{X}}_o|
    \bigl(\log n+\log k+\log l+c_F\bigr)
    +\log^*|\widehat{\mathcal{X}}_o|.
  \label{eq:cost-outlier}
\end{align}
Không cần thuộc từng dòng. Điều cần hiểu là: chi phí tăng theo số phần tử khác 0 và theo kích
thước các trục; do đó tăng $d_k,d_l,d_s$ làm mô hình đắt hơn, còn thêm outlier cũng làm mô hình
đắt hơn. Ký hiệu $n$ trong chi phí outlier không được định nghĩa lại ngay tại công thức; paper dùng
nó như độ dài stream/tổng số phần tử. Đây là một điểm chưa rõ của bài, nên tài liệu giữ nguyên
công thức và ghi chú thay vì tự gán nghĩa chắc chắn.

\subsection{Chi phí mã hóa dữ liệu}

Phần còn lại là chi phí mô tả dữ liệu dưới mô hình. Paper dùng mã Huffman và mô hình nhiễu Gauss
với trung bình $\mu$, phương sai $\sigma^2$:
\begin{equation}
  \bigl\langle\mathcal{X}\mid\mathcal{F}\bigr\rangle
  =\sum_{x\in\mathcal{X}}
   -\log_2 p_{\mu,\sigma}
   \bigl(x-\widehat{x}_d-\widehat{x}_s-\widehat{x}_o\bigr).
  \label{eq:data-cost}
\end{equation}
Đây là Eq. không số ngay sau Eq.~(10) trong paper. Trực giác: nếu residual
$x-\widehat{x}_d-\widehat{x}_s-\widehat{x}_o$ nhỏ và phù hợp với nhiễu dự kiến, nó tốn ít bit;
nếu residual lớn bất thường, nó tốn nhiều bit. Vì vậy mô hình vừa phải khớp, vừa không được biến
mọi dao động nhỏ thành tham số.

\begin{plainwords}
MDL không hỏi ``mô hình nào khớp nhất?''. Nó hỏi ``mô hình nào cho bản mô tả ngắn nhất của cả mô
hình lẫn dữ liệu?''. Một mô hình quá to có thể khớp hoàn hảo nhưng bản mô tả dài; một mô hình quá
nhỏ có bản mô tả ngắn nhưng phần residual dài. Điểm cân bằng chính là cấu trúc tối ưu.
\end{plainwords}

\subsection{Outlier và số chiều được chọn thế nào?}

Với outlier, paper thử đặt từng phần tử residual về 0 và giữ lại nếu tổng chi phí giảm. Đây là
``sparsification'' có tiêu chuẩn, không phải ngưỡng cảm tính.

Với số chiều, Algorithm~4 (\textsc{RankUpdate}) thử các cấu hình chỉ thay đổi một rank một đơn vị:
$(d_k-1,d_l,d_s)$, $(d_k+1,d_l,d_s)$, $(d_k,d_l-1,d_s)$, $(d_k,d_l+1,d_s)$,
$(d_k,d_l,d_s-1)$, $(d_k,d_l,d_s+1)$. Mỗi cấu hình được ước lượng lại, rồi giữ cấu hình có tổng
chi phí nhỏ nhất. Trong bản HTML/PDF của paper, danh sách candidate in ra lặp
$(d_k,d_l+1,d_s)$ và thiếu $(d_k+1,d_l,d_s)$; đây là lỗi trình bày có khả năng cao. Tài liệu dùng
danh sách đầy đủ theo ý tưởng ``mỗi lần đổi một rank'' và ghi chú rõ.

\subsection{ModelUpdate: thêm, giữ hay chuyển regime?}

Algorithm~3 so sánh tổng chi phí khi thêm candidate $\Theta'$ vào bộ mô hình $\mathcal{F}$ với
tổng chi phí khi không thêm. Nếu thêm làm chi phí giảm, candidate được chấp nhận; nếu không, mô
hình cũ được giữ. Sau khi chấp nhận, RankUpdate có thể điều chỉnh $d_k,d_l,d_s$.

\begin{center}
\resizebox{\textwidth}{!}{%
\begin{tikzpicture}[
  >=Latex,
  node distance=6mm and 6mm,
  stage/.style={draw=DTNavy, rounded corners=1.5pt, fill=DTMist,
    text width=3.2cm, minimum height=11mm, align=center, font=\sffamily\small},
  decision/.style={diamond, draw=DTClaret, fill=DTWarm, aspect=2.0,
    text width=2.2cm, align=center, font=\sffamily\small},
  arrow/.style={->, thick, DTTeal}
]
  \node[stage] (x) {cửa sổ mới\\$\mathcal{X}^{c}$};
  \node[stage, right=of x] (est) {ModelEstimation\\ứng viên $\Theta'$};
  \node[decision, right=of est] (cost) {chi phí\\giảm?};
  \node[stage, right=of cost] (add) {thêm $\Theta'$\\và RankUpdate};
  \node[stage, below=7mm of cost] (keep) {giữ mô hình cũ};
  \node[stage, below=7mm of add] (fc) {Forecast\\$L_f$ bước};
  \draw[arrow] (x) -- (est);
  \draw[arrow] (est) -- (cost);
  \draw[arrow] (cost) -- node[above,font=\small]{có} (add);
  \draw[arrow] (cost) -- node[right,font=\small]{không} (keep);
  \draw[arrow] (add) -- (fc);
  \draw[arrow, dashed, DTNavy] (fc.south) -- ++(0,-0.5) -| (x.south);
\end{tikzpicture}%
}
\end{center}

\subsection{Toàn bộ thuật toán streaming}

Algorithm~1 khởi tạo trên cửa sổ đầu tiên, đặt $\mathcal{F}=\{\Theta\}$, rồi với mỗi thời điểm
$t_c$:
\begin{enumerate}
  \item lấy cửa sổ $\mathcal{X}^{c}$;
  \item chạy ModelEstimation để có candidate $\Theta'$;
  \item chạy ModelUpdate để quyết định thêm/giữ và cập nhật rank;
  \item dự báo $L_f$ bước bằng $\mathcal{F}$;
  \item chuyển sang thời điểm kế tiếp.
\end{enumerate}

\begin{keypoint}
Chi phí mỗi bước không phụ thuộc vào tổng chiều dài stream vì bước 1 luôn lấy đúng $L_c$ lát cắt.
Tập $\mathcal{F}$ có thể tăng khi có regime mới, và RankUpdate tạo các ``spike'' thời gian, nhưng
không có vòng lặp trên toàn bộ lịch sử. Đây là ý nghĩa chính xác của scalability trong paper.
\end{keypoint}

\subsection{Forecast và những gì paper không nói}

Algorithm~1 gọi $\widehat{\mathcal{X}}_{t_c+L_f}=\mathrm{Forecast}(\mathcal{F})$ nhưng không trình
bày chi tiết công thức Forecast trong các mục chính. Vì vậy, tài liệu này chỉ nói điều chắc chắn:
dự báo dùng mô hình đang hoạt động trong $\mathcal{F}$, và trend có thể được ngoại suy từ hệ
reaction--diffusion. Ta \emph{không} tự bịa một công thức chi tiết cho seasonal/outlier. Nếu cần
tái hiện chính xác, phải đọc code gốc hoặc appendix/implementation của tác giả.

\subsection{``Parameter-free'' nghĩa là gì?}

Paper nói D-Tracker không cần chỉnh tham số. Cần đọc đúng ngữ cảnh: người dùng không phải chọn
tay số nhóm $d_k,d_l,d_s$ hay thời điểm đổi regime; các đại lượng đó được chọn bằng MDL. Tuy nhiên
vẫn có các lựa chọn kỹ thuật như bộ giải ODE, số vòng lặp ALS, dung sai hội tụ, và trong thí nghiệm
paper tìm kiếm khởi tạo $d_k,d_l\in[2,4]$, $d_s\in[0,4]$. ``Không hyperparameter'' không có nghĩa
``không có bất kỳ tham số số học nào''.

\begin{commonmistake}
Đừng nói ``MDL luôn tìm mô hình đúng''. MDL tối ưu một tiêu chuẩn mã hóa; nếu giả định nhiễu sai,
hoặc dữ liệu thật sự phức tạp hơn họ mô hình, mô hình tối ưu theo MDL vẫn có thể dự báo kém. Nó
cũng không thay thế kiểm định nhân quả cho các mũi tên diffusion.
\end{commonmistake}

\begin{selfcheck}
\begin{enumerate}
  \item Một mô hình tăng $d_k$ sẽ làm thành phần nào của MDL tăng, và có thể làm thành phần nào giảm?
  \item Vì sao outlier không được chọn bằng ngưỡng cố định?
  \item Nếu candidate model không làm giảm tổng chi phí, ModelUpdate làm gì?
  \item Giải thích trong một câu vì sao thời gian mỗi bước không tăng theo chiều dài stream.
\end{enumerate}
\end{selfcheck}
